\section{Token Template of Tool Calling}
\label{sec:tool-call-token-template}

There are three components in the token structure for tool-calling:
\begin{itemize}
    \item \textbf{Tool declaration message}: defines the list of available tools and the schema of the arguments;
    \item \textbf{Tool invoking section in assistant message}: encodes the model's request to invoke tools; 
    \item \textbf{Tool result message}: encapsulates the invoked tool's execution result.
\end{itemize}


\lstdefinestyle{mystyle}{
    keywordstyle=\footnotesize\ttfamily,
    basicstyle=\footnotesize\ttfamily,
    breakatwhitespace=false,         
    breaklines=true,                 
    captionpos=t,                    
    keepspaces=true,                 
    showspaces=false,                
    showstringspaces=false,
    showtabs=false,                  
    tabsize=2
}
\lstset{style=mystyle}


The raw tokens of the tool declaration message are formatted as follows:

\begin{center}
\framebox{%
\begin{minipage}{8cm}
\ttfamily\footnotesize
{\color{blue}<|im\_begin|>}\\
tool\_declare\\
{\color{blue}<|im\_middle|>}\\
\# Tools\\
\\
{\color{darkgreen}\{\{ tool declaration content \}\}}\\
{\color{blue}<|im\_end|>}
\end{minipage}
}
\end{center}

The blue highlighted marks represent special tokens, and the green part, quoted by brackets, is the tool declaration content.
We use TypeScript to express the tool declaration content, since TypeScript is a concise language with a comprehensive type system, able to express the types and constraints of tool parameters with brief text. The code~\ref{lst:tool-declare-json} shows
an example for two simple tools in JSON format compatible with OpenAI's chat completion API, as a comparison, the same tools defined in TypeScript (listed in Code~\ref{lst:tool-declare-ts}) is much shorter.
To improve compatibility, part of our training data also uses JSON as the tool declaration language, so that 3rd-party frameworks need not additional development to support our tool calling scheme.

\begin{lstlisting}[language=Python, caption=Tool definition with JSON in OpenAI compatible API, label={lst:tool-declare-json}]
[{
  "type": "function",
  "function": {
    "name": "get_weather",
    "description": "Get weather for a location and date",
    "parameters": {
      "type": "object",
      "properties": {
        "location": {
          "type": "string",
          "description": "City and country e.g. Beijing, China"
        },
        "date": {
          "type": "string",
          "description": "Date to query, format in `%Y-%m-%d'"
        }
      },
      "required": [
        "location"
      ]
    }
  }
},
{
  "type": "function",
  "function": {
    "name": "Calculator",
    "description": "Simple calculator",
    "parameters": {
      "properties": {
        "expr": {
          "type": "string",
          "description": "Arithmetic expression in javascript"
        }
      },
      "type": "object"
    }
  }
}]
\end{lstlisting}

\begin{lstlisting}[caption=Tool definition in TypeScript, label={lst:tool-declare-ts}]
namespace functions { 
// Get weather for a location and date
type get_weather = (_: { 
  // City and country e.g. Beijing, China 
  location: string, 
  // Date to query, format in `%Y-%m-%d' 
  date?: string 
}) => any; 
// Simple calculator 
type Calculator = (_: { 
  // Arithmetic expression in javascript 
  expr?: string 
}) => any; 
}
\end{lstlisting}


The token template of the tool invoking section in the model's response messages is listed as follows:
\begin{center}
\framebox{%
\begin{minipage}{14cm}
\ttfamily\footnotesize
{\color{blue}<|tool\_call\_section\_begin|>}\\
{\color{blue}<|tool\_call\_begin|>}\\
{\color{gray}  // call\_id part}\\
functions.{\color{darkgreen}\{\{tool name\}\}}:{\color{darkgreen}\{\{counter\}\}}\\
{\color{blue}<|tool\_arguments\_begin|>}\\
{\color{darkgreen}\{\{ json serialized call arguments \}\}}\\
{\color{blue}<|tool\_call\_end|>}\\\
{\color{blue}<|tool\_call\_begin|>}\\
{\color{gray}  // more tool calls}\\
{\color{blue}<|tool\_call\_end|>}\\
{\color{blue}<|tool\_call\_section\_end|>}
\end{minipage}
}
\end{center}

As shown in the template, we support parallel tool calling by placing multiple tool calls in a single response turn. 
Each tool call has a unique call id, formatted as \texttt{functions.\{tool-name\}:\{counter\}}, 
where \texttt{tool-name} is the name of the tool, and \texttt{counter} is an auto-increasing counter of all tool calls starting from 0 in the dialog.

During inference, the model may occasionally generate unexpected tokens, leading to format errors when parsing a tool call. To solve this issue, we developed a constrained decoding module named \emph{enforcer}, inspired by {lm-format-enforcer}\footnote{\url{https://github.com/noamgat/lm-format-enforcer}}. When a \texttt{<tool\_call\_section\_begin|>} token is generated, it ensures that the upcoming tool-related tokens follow the predefined template, and the JSON argument string follows the declared schema.

The tool result message is simply a text message encoded with the tool's call id and the corresponding results.
\begin{center}
\framebox{%
\begin{minipage}{8cm}
\ttfamily\footnotesize
{\color{blue}<|im\_begin|>}\\
tool\\
{\color{blue}<|im\_middle|>}\\
\#\# Results of {\color{darkgreen}\{\{call\_id\}\} }\\
{\color{darkgreen}\{\{ execution result content \}\}}\\
{\color{blue}<|im\_end|>}
\end{minipage}
}
\end{center}

